\section{Translation Symmetry and Energy--Momentum}
\label{sec:translation_energy_momentum}

A field theory that does not distinguish one spacetime point from another is invariant under translations. Noether's theorem then produces the energy-momentum tensor. Its time component describes energy density and energy flow, while its spatial components describe momentum density and stress.

\subsection{Infinitesimal translations}

Consider

\[
x'^{\mu}
=
x^{\mu}
+
\varepsilon a^{\mu},
\]

where \(a^{\mu}\) is constant. For scalar fields,

\[
\Psi_a=0.
\]

The characteristic is

\[
Q_a
=
-a^{\nu}\partial_{\nu}\phi_a.
\]

If the Lagrangian has no explicit coordinate dependence,

\[
\frac{\partial\mathcal{L}}
{\partial x^{\nu}}
=
0,
\]

then translations are symmetries.

\subsection{Noether current for translations}

The general current is

\[
j^{\mu}
=
\pi_a^{\mu}Q_a
+
\mathcal{L}\xi^{\mu}.
\]

Here

\[
\xi^{\mu}=a^{\mu},
\]

so

\[
j^{\mu}
=
-a^{\nu}
\pi_a^{\mu}\partial_{\nu}\phi_a
+
a^{\mu}\mathcal{L}.
\]

Factor out \(a^{\nu}\):

\[
j^{\mu}
=
-a^{\nu}
\left[
\pi_a^{\mu}\partial_{\nu}\phi_a
-
\delta^{\mu}{}_{\nu}\mathcal{L}
\right].
\]

Define the canonical energy-momentum tensor

\[
\boxed{
T^{\mu}{}_{\nu}
=
\frac{\partial\mathcal{L}}
{\partial(\partial_{\mu}\phi_a)}
\partial_{\nu}\phi_a
-
\delta^{\mu}{}_{\nu}\mathcal{L}
}.
\]

Then

\[
j^{\mu}
=
-a^{\nu}T^{\mu}{}_{\nu}.
\]

Because the constants \(a^{\nu}\) are independent, Noether conservation gives

\[
\boxed{
\partial_{\mu}T^{\mu}{}_{\nu}=0
}.
\]

There is one conservation law for each translation direction.

\subsection{Direct divergence check}

Differentiate:

\[
\partial_{\mu}T^{\mu}{}_{\nu}
=
\partial_{\mu}
\left(
\pi_a^{\mu}\partial_{\nu}\phi_a
\right)
-
\partial_{\nu}\mathcal{L}.
\]

Expand the first term:

\[
\partial_{\mu}T^{\mu}{}_{\nu}
=
\left(
\partial_{\mu}\pi_a^{\mu}
\right)
\partial_{\nu}\phi_a
+
\pi_a^{\mu}
\partial_{\mu}\partial_{\nu}\phi_a
-
\partial_{\nu}\mathcal{L}.
\]

If there is no explicit coordinate dependence,

\[
\partial_{\nu}\mathcal{L}
=
\frac{\partial\mathcal{L}}
{\partial\phi_a}
\partial_{\nu}\phi_a
+
\pi_a^{\mu}
\partial_{\nu}\partial_{\mu}\phi_a.
\]

Commuting partial derivatives, the last terms cancel:

\[
\partial_{\mu}T^{\mu}{}_{\nu}
=
\left[
\partial_{\mu}\pi_a^{\mu}
-
\frac{\partial\mathcal{L}}
{\partial\phi_a}
\right]
\partial_{\nu}\phi_a.
\]

Therefore

\[
\boxed{
\partial_{\mu}T^{\mu}{}_{\nu}
=
-\mathcal{E}_a
\partial_{\nu}\phi_a
}.
\]

On shell,

\[
\partial_{\mu}T^{\mu}{}_{\nu}=0.
\]

\subsection{Explicit coordinate dependence}

If the Lagrangian depends explicitly on \(x\), then on shell

\[
\boxed{
\partial_{\mu}T^{\mu}{}_{\nu}
=
-
\left.
\frac{\partial\mathcal{L}}
{\partial x^{\nu}}
\right|_{\mathrm{explicit}}
}.
\]

An externally prescribed inhomogeneous background can therefore exchange energy or momentum with the field.

\subsection{Real scalar field}

For

\[
\mathcal{L}
=
\frac{1}{2}
\partial_{\mu}\phi
\partial^{\mu}\phi
-
V(\phi),
\]

we have

\[
\pi^{\mu}
=
\partial^{\mu}\phi.
\]

Raising the second tensor index gives

\[
\boxed{
T^{\mu\nu}
=
\partial^{\mu}\phi
\partial^{\nu}\phi
-
\eta^{\mu\nu}\mathcal{L}
}.
\]

For a scalar field, this tensor is already symmetric:

\[
T^{\mu\nu}=T^{\nu\mu}.
\]

\subsection{Energy density}

Using the metric signature

\[
\eta_{\mu\nu}
=
\operatorname{diag}(1,-1,-1,-1),
\]

the Lagrangian density is

\[
\mathcal{L}
=
\frac{1}{2c^{2}}
\phi_t^{2}
-
\frac{1}{2}
|\nabla\phi|^{2}
-
V(\phi).
\]

The energy density is

\[
T^{00}
=
\partial^{0}\phi
\partial^{0}\phi
-
\eta^{00}\mathcal{L}.
\]

Since

\[
\partial^{0}\phi
=
\frac{1}{c}\phi_t,
\]

we obtain

\[
\boxed{
T^{00}
=
\frac{1}{2c^{2}}
\phi_t^{2}
+
\frac{1}{2}
|\nabla\phi|^{2}
+
V(\phi)
}.
\]

This agrees with the scalar-field energy density derived directly from the Hamiltonian.

\subsection{Energy flux}

For a spatial index \(i\),

\[
T^{i0}
=
\partial^{i}\phi
\partial^{0}\phi.
\]

Because

\[
\partial^{i}\phi
=
-\partial_i\phi,
\]

we find

\[
T^{i0}
=
-\frac{1}{c}
\phi_t\partial_i\phi.
\]

Define the energy-flux vector

\[
\boxed{
\mathbf{S}
=
c
\left(
T^{10},
T^{20},
T^{30}
\right)
=
-\phi_t\nabla\phi
}.
\]

The \(\nu=0\) conservation equation is

\[
\partial_{\mu}T^{\mu0}=0.
\]

Multiplying by \(c\),

\[
\boxed{
\frac{\partial T^{00}}{\partial t}
+
\nabla\cdot\mathbf{S}
=
0
}.
\]

This is the local energy balance law.

\subsection{Momentum density and total four-momentum}

The total four-momentum is

\[
\boxed{
P^{\nu}
=
\frac{1}{c}
\int_{\mathbb{R}^{3}}
T^{0\nu}
\,\mathrm{d}^{3}x
}.
\]

Its time component is

\[
P^{0}
=
\frac{E}{c},
\]

where

\[
E
=
\int
T^{00}
\,\mathrm{d}^{3}x.
\]

The spatial momentum density is

\[
\boxed{
\mathbf{g}
=
\frac{1}{c}
\left(
T^{01},
T^{02},
T^{03}
\right)
}.
\]

For the scalar field,

\[
T^{0i}
=
-\frac{1}{c}
\phi_t\partial_i\phi,
\]

so

\[
\boxed{
\mathbf{g}
=
-\frac{1}{c^{2}}
\phi_t\nabla\phi
}.
\]

\subsection{Plane-wave check}

Consider a massless plane wave in one dimension:

\[
\phi(x,t)
=
A\cos(kx-\omega t),
\qquad
\omega=ck.
\]

Then

\[
\phi_t
=
A\omega\sin(kx-\omega t),
\]

and

\[
\phi_x
=
-Ak\sin(kx-\omega t).
\]

The energy flux is

\[
S_x
=
-\phi_t\phi_x
=
A^{2}\omega k
\sin^{2}(kx-\omega t).
\]

Using \(\omega=ck\),

\[
S_x
=
cA^{2}k^{2}
\sin^{2}(kx-\omega t).
\]

The energy density is

\[
T^{00}
=
\frac{1}{2c^{2}}\phi_t^{2}
+
\frac{1}{2}\phi_x^{2}
=
A^{2}k^{2}
\sin^{2}(kx-\omega t).
\]

Therefore

\[
\boxed{
S_x=cT^{00}
}.
\]

The energy propagates in the positive \(x\) direction at speed \(c\), as expected.

\subsection{Canonical and symmetric tensors}

For scalar fields, the canonical tensor is symmetric. For fields carrying spin, the canonical tensor need not satisfy

\[
T^{\mu\nu}=T^{\nu\mu}.
\]

It can often be improved by adding a divergence term to obtain a symmetric tensor with the same total four-momentum. This issue becomes important for angular momentum and for coupling fields to gravity.

\subsection{Common mistakes}

Typical errors include:

\begin{enumerate}
    \item forgetting the term \(-\delta^{\mu}{}_{\nu}\mathcal{L}\);
    \item confusing \(T^{0i}\) with \(T^{i0}\) when the tensor is not symmetric;
    \item identifying \(T^{00}\) with the Lagrangian density;
    \item losing the minus sign in \(\partial^{i}=-\partial_i\);
    \item claiming momentum conservation in the presence of an explicitly position-dependent external source;
    \item omitting the boundary flux when deriving conservation of total energy.
\end{enumerate}

\subsection{What has been established}

Translation symmetry produces the canonical energy-momentum tensor

\[
T^{\mu}{}_{\nu}
=
\pi_a^{\mu}\partial_{\nu}\phi_a
-
\delta^{\mu}{}_{\nu}\mathcal{L}.
\]

For a translation-invariant theory,

\[
\partial_{\mu}T^{\mu}{}_{\nu}=0.
\]

For the real scalar field, \(T^{00}\) is the positive energy density and \(cT^{i0}\) is the energy flux. The next section uses rotational and Lorentz symmetry to construct angular-momentum currents.
